\section{Related work}
Persistent queues can make a busy link feel slow even when its throughput
looks healthy. Gettys and Nichols describe this bufferbloat problem and why
large buffers can hurt interactive traffic~\cite{bufferbloat}. Floyd and
Jacobson's RED work established the value of responding to congestion before
a queue overflows~\cite{red}. Nichols and Jacobson's CoDel instead uses packet
sojourn time to distinguish a standing queue from a transient burst~\cite{codel}.
These studies motivate our latency measurements and unloaded baselines.
We evaluate the deployed FQ-CoDel combination, rather than claiming to isolate
CoDel's contribution from flow scheduling.

Buffer size and traffic aggregation also matter. Appenzeller, Keslassy and
McKeown study how flow multiplexing changes router-buffer requirements~\cite{buffers}.
Their backbone assumptions do not map directly to our small testbed, but they
explain why the number of competing flows belongs in the design. We plan one
and eight background flows while holding aggregate offered demand fixed.

CAKE is the closest prior system to our experiment. H{\o}iland-J{\o}rgensen,
T{\"a}ht and Morton combine shaping, queue management and traffic isolation
for home gateways~\cite{cake}. We compare its best-effort mode with a bounded
classification rule for our foreground replies. The classification budget
matters because a priority policy that starves bulk traffic would fail our
goodput criterion. Our workload and repeatability checks are a small feasibility
study of these existing mechanisms, not a new queue-management algorithm.

DCTCP examines small foreground work competing with large transfers in data
centers~\cite{dctcp}. It changes endpoint congestion control using ECN, whereas
our main treatments change the Linux router. Dean and Barroso explain why
occasional slow responses can dominate a distributed service's user-visible
latency~\cite{tail}. Together, these papers support looking beyond averages,
but neither makes our 512-byte request stream a validated model of a complete
production service.

Finally, Flent shows why repeated network experiments need coordinated traffic,
metadata and reproducible analysis~\cite{flent}. Our collector follows that
measurement concern with retained attempts, hashes and explicit request
accounting. We use a custom foreground program and iperf3 rather than claiming
that our current pipeline runs Flent. Its methodology remains a useful
independent cross-check for the final study.
